\documentclass{article}

\usepackage[margin=1in]{geometry}
\usepackage{tabularx}
\usepackage{booktabs}
\usepackage[normalem]{ulem}

\title{Reflection and Traceability Report on BSS}

\author{Xinlu Yan}

\date{April 21, 2026}

\input{../Comments.text}
\input{../Common.text}

\begin{document}

\maketitle

\section{Changes in Response to Feedback}

All feedback is recorded on the
\href{https://github.com/Xinlu-Y/CAS741/issues}{GitHub issue tracker}; each
item is its own issue, and the response is either a comment, a linked
commit, or a tracking issue listing the resolving commits. Only
representative items are highlighted below.

\subsection{SRS \sout{and Hazard Analysis}}

The overall structure of my SRS stayed mostly the same after feedback.
That was partly because I had useful advice early on from Dr.~Smith and
from the course material, so the document began with a reasonable
structure.

Most of the feedback on the SRS was relatively minor. It mainly pointed
out mistakes, places where I had misunderstood the purpose of a section,
and places where the theory or its explanation needed to be clearer. The
largest changes were in the theoretical models, the instance model, and
the input constraints, which became much more explicit by the end.

The presentation feedback was especially helpful early in the term,
while I was still figuring out how to write the SRS. Getting that advice
early made it easier to continue the document and probably reduced the
number of larger issues in later rounds of feedback.

I did not receive meaningful feedback on the Hazard Analysis, which is
why it has been struck from the subsection title. Since BSS is only a
numerical simulation and does not interact with physical hardware, there
was not much hazard-related discussion.

\subsection{Design and Design Documentation}

There was no separate software design document for this project. Most
design-related feedback was absorbed into the SRS and later revisions to
the modelling assumptions, instance model, and validation strategy.

\subsection{VnV Plan \sout{and Report}}

The VnV~Plan received peer review and comments from Dr.~Smith. Most of the
feedback was relatively minor, including wording changes, numbering
fixes, cleanup, and small corrections throughout the document. The main
substantive comments focused on making the testing story clearer:
distinguishing validation from verification more carefully, placing
tests in the right categories, and justifying numerical tolerances and
oracle choices more explicitly. I revised the plan accordingly to make
those distinctions and justifications clearer.

I did not produce a VnV~Report for this project, which is why
\textit{Report} has been struck from the subsection title.

\section{Design Iteration (LO11 (PrototypeIterate))}

The design changed mostly through review, Drasil encoding, and VnV
planning, not through building a large prototype. After the first SRS
draft, I ended up restructuring the theoretical models, rewriting the
instance model into a cleaner ODE form, and making several constraints
more explicit.

The main change was that some assumptions stopped being just background
statements and became real design constraints. The center-of-mass frame
needed zero-total-momentum initial conditions to be checkable, and the
nonzero-separation assumption had to be treated as a real precondition
for the force model. Drasil also forced some changes in how I wrote the
design, since some mathematical expressions did not fit cleanly into the
available representations. The VnV plan pushed this a bit further by
making me separate input handling, solver behaviour, verification, and
validation more clearly.

\section{Design Decisions (LO12)}

The biggest design decision was to keep the problem narrow. I focused on
a two-body binary star system in 2-D under Newtonian gravity, instead of
trying to support a more general astrophysical setting. That meant
leaving out relativity, third-body effects, collisions, and similar
complications so the project would stay manageable.

Another decision was to turn some modelling assumptions into explicit
constraints. The center-of-mass frame, planar motion, and nonzero
separation all affect the inputs, the equations, and what counts as a
valid result, so I did not want them to remain as vague background
assumptions.

I also chose to keep the SRS focused on the problem and the model,
rather than forcing it too early into an implementation-oriented form.
At the same time, I still had to work within Drasil's current limits,
especially for vectors, function representations, and differential
equation expressions. In practice, that meant some formulations had to
be rewritten into forms that Drasil could handle better.

Finally, I treated the IVP solver as external to the scientific model.
The software is responsible for checking inputs and setting up the
problem correctly, while the solver is responsible for attempting to
solve the ODE system.

\section{Economic Considerations (LO23)}

BSS is not a commercial product. Its target audience is undergraduate
students taking a first course in classical mechanics, scientific
computing, or numerical methods, where a small, transparent, well-documented
two-body simulator is more useful as a teaching aid than a production
astrophysics package would be. That market is already well served by
free open-source tools, so BSS would not attract paying users on its
own. What matters economically is its role as a Drasil case study:
producing a generated scientific computing program from a documented
SRS for a non-trivial physics problem strengthens the argument that
Drasil's document-first approach can scale. To attract users to BSS
specifically, the most promising route is releasing it on GitHub with
the User Guide as the entry point and a few worked example inputs
(circular, eccentric, unequal-mass), so any undergraduate physics or
computational science course can adopt it cheaply.

\section{Reflection on Project Management (LO24)}

\subsection{How Does Your Project Management Compare to Your Development Plan}


As the sole member of this project, the team-oriented elements of the Development Plan,
such as the meeting plan, communication plan, and role assignments, were not directly
applicable. Issue tracking, however, remained highly relevant: I used GitHub issues
throughout the project, turning each piece of feedback into an issue that was then
resolved either by a referenced commit or by linking it to a tracking issue.

The technology stack ultimately matched the original plan, consisting of Python, LaTeX,
and Drasil for generating the SRS. Coordination with Zhuo Zhang and Dr. Smith was also
smooth. The course community was active and collaborative on GitHub, which made online
discussions efficient and productive.

\subsection{What Went Well?}

What went well for me in this project was my ability to move beyond just learning the toolchain and
begin understanding the deeper ideas behind Drasil. Once I became comfortable with the technologies
and workflow, I was able to focus much more on how the system organizes knowledge through chunk
types and how that structure supports a stronger SRS. Building the BSS example helped me develop
that understanding in a concrete way, since it required me to think carefully about theory models,
derivations, assumptions, and consistency across the document. I was especially happy that this
understanding transferred into my code generation work as well. Rather than treating the example and
the generator as separate tasks, I was able to apply the same knowledge-first perspective to fix
several code generation issues, including problems with multi-value outputs and ODE solution
handling across languages. That transfer of knowledge was one of the most successful parts of the
project for me, because it showed real growth in both my understanding of Drasil and my ability to
contribute to it in a meaningful way.

\subsection{What Went Wrong?}

Although the project did not encounter any major breakdowns in process or tooling, a few challenges did affect progress. The main difficulty was the Drasil encoding. The most significant issue was the
gap between the level of precision I wanted in the SRS and mathematical development and the form that Drasil currently handles most naturally for code
generation. In particular, areas such as function definitions, vectors, and manually written code snippets required extra effort and did not always fit smoothly into the existing workflow. Some constructs I wanted
(a vector gravitational law with an explicit unit vector, traceability
links to non-UID property statements) had no clean Drasil counterpart
and I had to work around them. The PDF output of the generated SRS
also had table-overflow issues with the traceability matrices.

\subsection{What Would you Do Differently Next Time?}

  Next time, I would keep the same document-driven development approach, but I would structure the project around the known limitations of Drasil much earlier.
  In particular, I would narrow the problem scope sooner and make earlier decisions about how mathematical objects such as functions and differential equation
  solutions should be represented within Drasil. Doing so would reduce rework and make the transition from SRS development to implementation smoother.
  One especially valuable improvement would be a guided tutorial for new Drasil authors, helping them set up a local development
  environment and build a small working example that can serve as a template for their own projects.

  More broadly, this project convinced me of the value of a principled, document-driven approach to software development. I would use this approach again in
  future work.

\end{document}
